\documentclass[version=final]{iacrcc}

\license{CC-by}

% Packages NOT included by iacrcc
\usepackage{booktabs}
\usepackage{array}
\usepackage{enumitem}
\usepackage{algorithm}
\usepackage{algpseudocode}
% cleveref is loaded via after-hyperref.sty (required by iacrcc)

% Observation is not defined by iacrcc.  It shares the theorem counter, and
% cleveref (loaded by the class after the preamble) derives its reference type
% from that counter, so \Cref would render "Theorem" for an observation.  All
% cross-references to the observation below are therefore written explicitly as
% "Observation~\ref{...}".
\newtheorem{observation}[theorem]{Observation}

% Math operators and commands from the paper
\DeclareMathOperator{\tr}{tr}
\DeclareMathOperator{\diag}{diag}
\DeclareMathOperator{\Id}{Id}
\DeclareMathOperator{\negl}{negl}
\newcommand{\norm}[1]{\lVert #1 \rVert}
\newcommand{\abs}[1]{\lvert #1 \rvert}
\newcommand{\frob}[1]{\norm{#1}_{\mathrm{F}}}
\newcommand{\opn}[1]{\norm{#1}_{\mathrm{op}}}
\newcommand{\dtv}{d_{\mathrm{TV}}}
\newcommand{\Dhat}{\widehat{D}}
\newcommand{\CC}{\mathrm{CC}}
\newcommand{\BP}{\mathrm{BP}}

\title[running={Uniform Mixing in Barrington Programs}]{Uniform Mixing in Barrington Branching Programs}
\addauthor[inst={1},email={0244552@up.edu.mx},surname={Zermeño}]{Ernest Darell Zermeño}
\addaffiliation[country={Mexico}]{Universidad Panamericana}

\begin{document}

\maketitle

\keywords[Barrington branching programs, spectral contraction, Fourier analysis on finite groups, symmetric group, indistinguishability obfuscation, point functions]{Barrington branching programs, spectral contraction, Fourier analysis on finite groups, symmetric group, indistinguishability obfuscation, point functions}

\begin{abstract}
We study the distribution $D_T$ over the symmetric group $S_5$ induced by evaluating
any fixed valid balanced transposition-template Barrington branching program for the
point function $\CC_T$ on uniformly random inconsistent inputs. We prove two
results for this template. First, $D_T = D_{T'}$ for all
$T, T' \in \{0,1\}^\ell$ (exact T-independence), via a combinatorial flip bijection
that requires no properties of the underlying group. Second,
$\dtv(D_T, \mathrm{Uniform}(S_5)) \leq \negl(\ell)$ (spectral contraction),
via Fourier analysis on $S_5$ using the Diaconis--Shahshahani Upper Bound Lemma.
The contraction proof proceeds by exact computation of Fourier norms for the standard
representation $(4,1)$ at the first two levels of the commutator tree, followed by
exhaustive verification that the operator norm drops below~$1$ at the third level
--- enabled by the discovery that the level-$2$ obstruction has a regular simplex
structure in $\mathbb{R}^4$ that the commutator necessarily breaks. We further
identify obstructions to extending this uniform contraction beyond $S_5$:
for $S_n$ with $n \geq 6$ a dimensional obstruction forces every
level-$2$ operator to have norm~$1$. At level~$3$, the maximum norm is~$1$
for $S_6$ and $S_7$, although some configurations contract; for $n \geq 10$
every configuration has norm~$1$ by a dimensional bound. The analysis presented
here does not decide the intermediate cases $n \in \{8,9\}$.
The two results combine into a quantitative statement about one exactly specified
experiment: the $q$-sample transcript of terminal outputs obtained from fresh,
independent, uniformly random inconsistent paths is exactly target-independent
and within $\min\{1, q\,\varepsilon(\ell)\}$ of uniform, where $\varepsilon(\ell)$
is our explicit single-sample bound (so the transcript remains negligibly close
to uniform whenever $q = q(\ell)$ is polynomially bounded), and both properties survive
target-independent post-processing. This gives an unconditional, quantitative
terminal-output guarantee for the fixed transposition-template Barrington layer.
\end{abstract}

\begin{textabstract}
We study the output distribution of a fixed valid balanced transposition-template Barrington branching program for a point function on uniformly random inconsistent inputs. We prove exact target-independence by a combinatorial flip bijection and negligible statistical distance from uniform on S_5 by finite-group Fourier analysis. A regular simplex describes the level-2 obstruction in the standard representation; the commutator structure breaks it at level 3. An exact integer certificate gives the uniform contraction constant 0.0355, and induction amplifies the contraction at higher levels. We also identify dimensional obstructions beyond S_5: the maximum level-3 norm is 1 for S_6 and S_7, although some configurations contract, and every level-3 configuration has norm 1 for n at least 10. The present analysis does not decide S_8 and S_9. Fresh independent terminal-output samples are exactly target-independent and remain negligibly close to uniform for polynomially many samples. These properties survive target-independent post-processing.
\end{textabstract}


%% ===================================================================
\section{Introduction}
\label{sec:intro}
%% ===================================================================

Indistinguishability obfuscation (iO) for general circuits is one of the most
powerful primitives in cryptography, implying a vast array of applications from
functional encryption to deniable encryption~\cite{GGH+13,SW14}. The first
candidate constructions~\cite{GGH+13} were based on matrix branching programs
over multilinear maps, randomized via Kilian's
technique~\cite{Kil88}, following the classical correspondence established by
Barrington's theorem~\cite{Bar89}: every $\mathrm{NC}^1$ circuit of depth $d$
can be simulated by a branching program of width~$5$ and length at most $4^d$
over the symmetric group $S_5$.

Obfuscation of branching programs, and the mixed-input attacks studied against
it~\cite{CVW18}, motivate the question we ask here: what distribution does
such a program produce when it is evaluated on ``incorrect'' inputs?
Specifically, for a point function $\CC_T : \{0,1\}^\ell
\to \{0,1\}$ that outputs $1$ only on $x = T$, we study the fixed balanced
transposition-template program $\BP_T$ described below: it produces a designated
permutation $\pi_{\mathrm{accept}} \in S_5$ on input $T$ and the identity $e$ on
all other inputs. When this program is evaluated on a random bit string
$b \in \{0,1\}^L$ that is \emph{inconsistent} (at least one variable receives
different values at different steps), the output is a random element of $S_5$
whose distribution we denote $D_T$.

Two properties of $D_T$ are natural to ask of this layer, and are the ones we
establish:

\begin{enumerate}[label=(\roman*)]
\item \textbf{T-independence:} $D_T = D_{T'}$ for all $T, T'$: the marginal law
  of the terminal output on a uniformly random inconsistent path is
  \emph{exactly} independent of the target.
\item \textbf{Uniform mixing:} $\dtv(D_T, \mathrm{Uniform}(S_5)) \leq \negl(\ell)$,
  so that marginal is statistically close to uniform on $S_5$.
\end{enumerate}

\paragraph{The experiment we study.}
Mixed-input evaluation---running a branching program on inputs that read a
variable inconsistently---is the setting in which the algebraic behaviour of the
$S_5$ layer is examined~\cite{CVW18}, and it is what motivates the question we
ask. We fix one exactly specified experiment and analyse it in full: the
$q$-sample \emph{terminal-output transcript}, in which each sample is the final
permutation produced by $\BP_T$ on a fresh, independent, uniformly random
inconsistent path (\Cref{def:terminal}). For the fixed balanced transposition
template we prove that this transcript is \emph{exactly} target-independent and
within $\min\{1, q\,\varepsilon(\ell)\}$ of uniform on $S_5^q$, where
$\varepsilon(\ell)$ is the explicit single-sample bound of \Cref{thm:main}(b);
in particular the transcript stays negligibly close to uniform whenever
$q = q(\ell)$ is polynomially bounded. Both properties are preserved by any
target-independent post-processing of the transcript (\Cref{cor:terminal}).

Uniform sampling conditioned on inconsistency is an analytical average-case
choice: it isolates the distributional behaviour of the fixed
transposition-template Barrington layer outside consistent Boolean evaluations.
We are not aware of an existing obfuscation security analysis that induces
exactly this sampling law. The contribution is therefore a mathematical result
about Barrington programs motivated by obfuscation.

\paragraph{A construction-level illustration.}
The GGH-style candidate of Garg, Gentry, Halevi, Raykova, Sahai, and
Waters~\cite[full version, Sections~4.1--4.2]{GGH+13}
illustrates the objects a security analysis must relate. Its public garbled
program contains encodings of randomized step matrices for the available input
choices, encoded endpoint vectors, a parallel dummy program, and public encoding
parameters. For an unfixed input bit both choices are available; evaluation
combines these encodings and tests the resulting expression. An adversary
receives this entire collection and can form expressions across evaluations.

Our theorem controls the terminal permutation of the plaintext template defined
here under fresh uniform inconsistent-path sampling. This supplies a model for
one underlying group-product distribution, without identifying our template
with the branching program of that candidate. A cryptographic application would
require a candidate-specific reduction relating the encoded view and its
permitted operations to this sampling experiment, and controlling information
beyond the terminal product. The terminal marginal alone does not establish
statistical obfuscation or the security of an iO construction.

\paragraph{Our contributions.}
We prove both properties for every fixed valid transposition template (\Cref{def:validtemplate}).

For T-independence, we give a short combinatorial proof (\Cref{thm:t-indep})
via a flip bijection $F_S$ that maps inconsistent paths for $T$ to
inconsistent paths for $T'$ while preserving the output permutation. The proof
uses no property of $S_5$ and extends to conjunctions with wildcards.

For uniform mixing, we develop a Fourier-analytic framework on $S_5$
(\Cref{sec:contraction}). Using the Diaconis--Shahshahani Upper Bound
Lemma, we reduce the problem to bounding the Fourier norms of the unconditional
distribution $P_T$ in each irreducible representation, then transfer the bound to
$D_T$ via a conditioning lemma (\Cref{lem:conditioning}). The standard
representation $(4,1)$ of dimension~$4$ is the bottleneck; we therefore focus the
exposition on it, while \Cref{prop:other_reps} and the level-$3$ certificate control
all remaining non-trivial irreducibles, as the Diaconis--Shahshahani bound requires. We establish contraction at the first
two levels of the commutator tree by exact symbolic computation
(\Cref{thm:level01,thm:level12}), discover that the level-$2$ obstruction
has a regular simplex structure (Observation~\ref{obs:simplex}), and prove that this
structure is necessarily broken at level~$3$ (\Cref{thm:level23}). An exact
integer Gram-form certificate over the full ordered universe then yields
operator-norm contraction with $\lambda \leq 0.0355$. An inductive argument
then gives doubly-exponential decay in the depth (\Cref{thm:induction}).

The mixing bound holds for powers of two $\ell \geq 8$; the sampling hypotheses
are collected in \Cref{rem:terminal-limits}. T-independence also extends to
conjunctions with a fixed wildcard set. \Cref{sec:extensions} presents evidence
for mixing in that setting and questions about other programs and finite groups.

\paragraph{Related work.}
Black-box obfuscation is impossible for general circuits, as shown by Barak,
Goldreich, Impagliazzo, Rudich, Sahai, Vadhan, and Yang~\cite{BGHSV01}; this
motivates the weaker notion of indistinguishability obfuscation.
The first iO candidate from branching programs~\cite{GGH+13} was followed
by constructions for specific function
classes: lockable obfuscation for point functions from
LWE~\cite{GKW17,WZ17}, conjunction obfuscation from noisy encodings
and LPN~\cite{BKMPRS18,BLMZ19}, and entropic Ring-LWE
approaches~\cite{BVWW16}.  The landmark result of Jain, Lin, and
Sahai~\cite{JLS21} achieved iO for general circuits from well-founded
assumptions, departing from the branching-program paradigm.  Recent
lattice-based candidates~\cite{BDGM20,WW21,HJL25} depart from the $S_5$
branching-program paradigm and rest on lattice assumptions, while impossibility
results~\cite{DJMMV25} constrain the space of viable assumptions. Mixed-input
evaluation --- running the program on inputs that read a variable
inconsistently --- is the setting in which the algebraic behaviour of a
branching-program layer is analysed~\cite{CVW18}, and is the setting of the
experiment we study.  Annihilation attacks~\cite{MSZ16} demonstrated
concrete vulnerabilities in early branching-program obfuscation candidates,
underscoring the importance of understanding the algebraic structure of
inconsistent evaluations.

\paragraph{Random subproducts.}
The unconditional law $P_T$ studied here is a \emph{random subproduct}: a fixed
word $g_1^{\varepsilon_1}\cdots g_L^{\varepsilon_L}$ in which each prescribed
group element is independently activated, $\varepsilon_i \sim \mathrm{Ber}(1/2)$.
Such products have been studied as lazy-transposition shuffles and, more
generally, under the name of group mixability~\cite{AH18,ABGV25}; those works ask
whether one can \emph{choose} a sequence---typically together with its activation
probabilities---whose product is exactly uniform, and show in particular that
symmetric groups are mixable. Our setting differs in both respects: every
activation probability is fixed to $1/2$, and the word is not free but is the
recursive expansion of a valid balanced Barrington commutator template. Exact
uniformity is then impossible for $P_T$, since each of its atoms has dyadic
probability while $1/\abs{S_5} = 1/120$ is not dyadic. We instead prove an
approximate-mixing bound, uniform over this constrained family, and transfer it
to the conditioned law $D_T$ (\Cref{lem:conditioning}).

Fourier-analytic methods for branching programs have a separate tradition
in the pseudorandomness literature~\cite{RSV13,LPV22}, but those
works use Fourier analysis on $\{0,1\}^n$ (the input space) to fool
Boolean outputs of read-once programs.  Our analysis instead uses Fourier
analysis on $S_5$ (the output group) to bound total variation of the
group-valued output distribution of multi-read programs --- a fundamentally
different analytic object requiring different tools.
Canetti, Chamon, Mucciolo, and Ruckenstein~\cite{CCMR24} recently proposed a ``local mixing''
framework for general-purpose obfuscation based on reversible circuits and
thermodynamic arguments; their mixing notion is structural (circuit-level),
whereas ours is distributional (output-group-level), making the two
approaches complementary. Our analysis applies the Diaconis--Shahshahani
framework~\cite{DS81} to the output group of the program.

\paragraph{Techniques.}
The proof of T-independence is purely combinatorial. The proof of spectral
contraction uses three different methods at three different scales:
exact symbolic algebra (levels~$0$--$2$), analytic structure and exact integer
certification (level~$3$), with independent exhaustive numerical cross-checks,
and analytic induction (levels~$4$+). The key discovery
that bridges levels~$2$ and~$3$ is the simplex structure of the
preserved directions: in the reduced level-$2$ census, the $270$ out of $405$
configurations with operator norm~$1$ each preserve one of exactly $5$ directions in
$\mathbb{R}^4$, forming a regular simplex with pairwise $\abs{\cos\theta} = 1/4$.
An algebraic proof shows that valid level-$3$ configurations cannot share a
preserved direction: within the stabilizer $\mathrm{Stab}(k) \cong S_4$,
the derived series terminates at depth~$3$, forcing the outer commutator to be
trivial. A separate analysis of $S_n$ for $n \geq 6$ shows that this mechanism
is specific to $S_5$: the level-$1$ Fourier transform preserves a subspace of
dimension $n-3$ in $\mathbb{R}^{n-1}$, and the intersection bound
$2(n-3)-(n-1) = n-5 \geq 1$ forces $\opn{\Dhat_2} = 1$ for every level-$2$
quadruple. At level~$3$, explicit norm-$1$ configurations rule out uniform
contraction for $S_6$ and $S_7$, whereas a further dimensional bound forces
every configuration to have norm~$1$ for $n \geq 10$. The analysis here does
not decide $n \in \{8,9\}$.


%% ===================================================================
\section{Preliminaries}
\label{sec:prelim}
%% ===================================================================

\subsection{Branching Programs and Barrington's Theorem}

\begin{definition}[Permutation Branching Program]
A \emph{permutation branching program} of width $w$ and length $L$ over a group
$G \leq S_w$ is a sequence of instructions
$\Pi = ((i_1, g_1^{(0)}, g_1^{(1)}), \ldots, (i_L, g_L^{(0)}, g_L^{(1)}))$,
where each $i_j \in [\ell]$ is a variable index and $g_j^{(0)}, g_j^{(1)} \in G$.
On input $x \in \{0,1\}^\ell$, the program evaluates
$\Pi(x) = \prod_{j=1}^{L} g_j^{(x_{i_j})}$.
The program $\alpha$-\emph{computes} a function $f$ if $\Pi(x) = \alpha$ when
$f(x) = 1$ and $\Pi(x) = e$ when $f(x) = 0$.
\end{definition}

\begin{theorem}[Barrington~\cite{Bar89}]
\label{thm:barrington}
Every Boolean function computed by a circuit of depth $d$ can be
$\alpha$-computed by a permutation branching program of width $5$ and
length at most $4^d$ over $S_5$, for some $\alpha \neq e$.
\end{theorem}

The construction uses the \emph{commutator trick}: given programs $P$ and $Q$
that $\alpha$-compute $f$ and $\beta$-compute $g$ respectively, the concatenation
$P \cdot Q \cdot P^{-1} \cdot Q^{-1}$ computes $f \wedge g$ via the group
commutator $[\alpha, \beta] = \alpha\beta\alpha^{-1}\beta^{-1}$. This works because
$S_5$ contains enough non-commuting elements to choose $\alpha,\beta$ and their
conjugates so that the required commutators remain non-identity at each gate.
Non-identity elements can of course commute; the Barrington construction depends
on these careful choices rather than on an arbitrary commutator being nontrivial.

\begin{definition}[Valid balanced transposition template]
\label{def:validtemplate}
A \emph{balanced transposition template} $\mathcal{T}$ of level $r$ assigns a
transposition of $S_5$ to each leaf of the complete binary tree of depth $r$,
together with a fixed orientation at each node. Its \emph{target} is defined
recursively: a level-$0$ configuration is a transposition, with target itself;
a level-$(r{+}1)$ configuration is an ordered pair $(A,B)$ of level-$r$
configurations, with target $\alpha_{(A,B)} = [\alpha_A, \alpha_B]$. The
template is \emph{valid} if every configuration in it---each sub-block and the
root---has non-identity target; equivalently, $\alpha_A \neq e$, $\alpha_B \neq e$
and $[\alpha_A,\alpha_B] \neq e$ at every node. We refer to the valid
configurations of a given level as valid \emph{pairs} (level~$1$),
\emph{quadruples} (level~$2$) and \emph{level-$3$ configurations}.

All point-function mixing bounds below are proved \emph{for every fixed valid
balanced transposition template}: a template $\mathcal{T}$ is fixed first (it is
what the comparison of two targets $T$ and $T'$ holds constant), and the bounds
then hold uniformly over all valid templates. The notation $D_T$ suppresses the
dependence on $\mathcal{T}$, which is fixed throughout. (Statements of wider
scope---T-independence for a common instruction template, conjunctions, and other
groups---are stated with their own hypotheses where they appear.)
\end{definition}

\paragraph{Non-vacuity.}
Valid templates exist at every level $r\geq3$.  On the points
$\{0,1,2,3,4\}$, the level-$3$ tree
\[
\mathcal{T}_3=
\Bigl[
  \bigl[[(01),(02)],[(01),(04)]\bigr],
  \bigl[[(01),(02)],[(01),(13)]\bigr]
\Bigr]
\]
is valid and has target $\rho=(01234)$.  Moreover,
$[(02431),(02314)]=\rho$, and both inputs are conjugate to~$\rho$.
Simultaneously conjugating all leaves of a valid template by a fixed group
element preserves transposition leaves and validity while conjugating every
internal target.  Thus two suitably conjugated copies of a level-$r$ template
with target $\rho$ form a valid level-$(r+1)$ template with the same target;
iteration proves the claim.

\begin{definition}[Point Function]
The point function $\CC_T : \{0,1\}^\ell \to \{0,1\}$ is defined by
$\CC_T(x) = 1$ if and only if $x = T$.
\end{definition}

For point functions, we restrict to the following fixed balanced
transposition-template construction. The template has $\ell$ leaves (one per bit
position), each comparing $x_i$ to $T_i$. Leaf $i$ outputs a fixed transposition
$\sigma_i$ when $x_i = T_i$ and the identity $e$ when $x_i \neq T_i$. The leaves
are combined by a fixed balanced binary tree of commutators, with a fixed
orientation at every internal node. When $\ell$ is a power of two, the tree has
depth $\log_2 \ell$, yielding a program of length $L = 4^{\log_2 \ell}=\ell^2$.
All statements below concern this template and the branch swaps induced by
changing $T$. They do not claim uniform mixing for arbitrary width-$5$
Barrington programs or for templates using only even permutations such as
$5$-cycles.


\subsection{\texorpdfstring{The Distribution $D_T$}{The Distribution D\_T}}

\begin{definition}[Inconsistent Path]
A bit string $b = (b_1, \ldots, b_L) \in \{0,1\}^L$ is \emph{inconsistent}
with respect to the program $\Pi$ if there exist steps $j, j'$ reading the
same variable ($i_j = i_{j'}$) with $b_j \neq b_{j'}$.
\end{definition}

\begin{definition}[$D_T$]
Let $\mathcal{I} \subset \{0,1\}^L$ be the set of inconsistent paths.
The instructions $g_j^{(\cdot)}$ are those of the program $\BP_T$ for the point
function $\CC_T$; the subscript $T$ records this dependence (by
\Cref{thm:t-indep}, $D_T$ is in fact independent of $T$).
The distribution $D_T$ over $S_5$ is defined by
\[
D_T(\pi) = \Pr_{b \sim \mathrm{Uniform}(\mathcal{I})}
\left[ \prod_{j=1}^{L} g_j^{(b_j)} = \pi \right].
\]
\end{definition}


\subsection{\texorpdfstring{Representation Theory of $S_5$}{Representation Theory of S5}}

The symmetric group $S_5$ has $7$ irreducible representations, labeled by
partitions of $5$:

\begin{center}
\begin{tabular}{lccl}
\toprule
Partition $\lambda$ & Dimension $d_\lambda$ & $\chi_\lambda(\text{transposition})$ & Name \\
\midrule
$(5)$       & 1 & 1 & Trivial \\
$(4,1)$     & 4 & 2 & Standard \\
$(3,2)$     & 5 & 1 & --- \\
$(3,1,1)$   & 6 & 0 & --- \\
$(2,2,1)$   & 5 & $-1$ & Sign $\otimes$ $(3,2)$ \\
$(2,1,1,1)$ & 4 & $-2$ & Sign $\otimes$ Standard \\
$(1^5)$     & 1 & $-1$ & Sign \\
\bottomrule
\end{tabular}
\end{center}

The dimensional check confirms: $1^2 + 4^2 + 5^2 + 6^2 + 5^2 + 4^2 + 1^2 = 120 = |S_5|$.

We call two transpositions $\sigma, \tau \in S_5$ a \emph{non-commuting pair}
if $[\sigma, \tau] \neq e$, equivalently if they share exactly one element.
There are $30$ such unordered pairs among the $10$ transpositions of $S_5$.

For a matrix $M \in \mathbb{R}^{d \times d}$, we write
$\frob{M}^2 = \tr(M^\top M) = \sum_{i,j} M_{ij}^2$ for the squared Frobenius
norm and $\opn{M} = \max_{\norm{v}=1} \norm{Mv}$ for the operator norm
(largest singular value).  The key inequality relating them is
$\opn{M}^2 \leq \frob{M}^2 \leq d \cdot \opn{M}^2$.

\begin{definition}[Standard Representation]
The standard representation $\rho_{(4,1)} : S_5 \to \mathrm{GL}(\mathbb{R}^4)$
acts on the hyperplane $V = \{v \in \mathbb{R}^5 : \sum_i v_i = 0\}$ by
permuting coordinates. We use the orthonormal basis
\begin{align*}
f_1 &= (e_0 - e_1)/\sqrt{2}, \quad
f_2 = (e_0 + e_1 - 2e_2)/\sqrt{6}, \\
f_3 &= (e_0 + e_1 + e_2 - 3e_3)/\sqrt{12}, \quad
f_4 = (e_0 + e_1 + e_2 + e_3 - 4e_4)/\sqrt{20},
\end{align*}
so that $\rho_{(4,1)}(\sigma)$ is orthogonal for all $\sigma \in S_5$.
\end{definition}


\subsection{The Diaconis--Shahshahani Upper Bound Lemma}

\begin{theorem}[Diaconis--Shahshahani~\cite{DS81}]
\label{thm:ds}
For any probability distribution $Q$ on a finite group $G$,
\[
4 \cdot \dtv(Q, U)^2 \leq \sum_{\rho \neq \mathrm{trivial}}
d_\rho \cdot \frob{\widehat{Q}(\rho)}^2,
\]
where the sum is over non-trivial irreducible representations $\rho$ of
dimension $d_\rho$, $\widehat{Q}(\rho) = \sum_{g \in G} Q(g) \rho(g)$ is the
Fourier transform, and $\frob{\cdot}$ denotes the Frobenius norm.
\end{theorem}

For $S_5$, the standard representation $(4,1)$ is the \emph{mixing
bottleneck}. Although $(2,1,1,1) = \mathrm{sign} \otimes (4,1)$ shares the
same normalized character ratio $\abs{\chi(\sigma)}/d_\rho = 1/2$, its leaf
projection $\frac{1}{2}(I - \rho_{(4,1)}(\sigma))$ has rank~$1$ (versus
rank~$3$ for $(4,1)$), causing it to contract at level~$1$ with operator
norm $\leq (1/2)^3 = 1/8$. The remaining representations satisfy
$\abs{\chi(\sigma)}/d_\rho \leq 1/5$ and likewise contract earlier
(\Cref{prop:other_reps}). We therefore focus the exposition on
$\frob{\Dhat(4,1)}^2$, while \Cref{prop:other_reps} and the level-$3$
certificate control every other non-trivial irreducible, as the
Diaconis--Shahshahani bound requires.


%% ===================================================================
\section{Exact T-Independence}
\label{sec:t-indep}
%% ===================================================================

We first show that the output distribution carries no information about the
target: for programs built from a common template, $D_T$ is literally
independent of~$T$. The argument is a short combinatorial bijection that uses no
property of~$S_5$, so we give it before the group-specific mixing analysis.

\begin{theorem}[T-Independence]
\label{thm:t-indep}
Let $\BP_T$ be the Barrington branching program for the point function $\CC_T$
over $\{0,1\}^\ell$, with all programs built from the same instruction template:
the same variable-read sequence and the same transposition at each leaf, with the
target $T$ only determining which branch receives that transposition. Then
$D_T = D_{T'}$ for all $T, T' \in \{0,1\}^\ell$.
\end{theorem}

\begin{proof}
Let $S = \{i : T_i \neq T'_i\}$ be the set of differing bit positions.
Define $F_S : \{0,1\}^L \to \{0,1\}^L$ by
\[
F_S(b)_j =
\begin{cases}
1 - b_j & \text{if } i_j \in S, \\
b_j & \text{otherwise},
\end{cases}
\]
where $i_j$ is the variable read at step $j$.

We verify three properties:

\textbf{(1) $F_S$ is an involution on $\{0,1\}^L$.}
Applying $F_S$ twice flips each bit in $S$-positions twice, returning to
the original string. Hence $F_S \circ F_S = \Id$, and in particular $F_S$
is a bijection on the full space $\{0,1\}^L$.

\textbf{(2) $F_S$ preserves inconsistency.}
$F_S$ flips \emph{all} readings of each variable $x_i$ with $i \in S$
simultaneously. If all readings of $x_i$ were equal (say, all $0$), after
flipping they are all $1$ --- still consistent. If some readings differed,
after flipping they still differ. Variables with $i \notin S$ are unaffected.
Therefore inconsistency of each variable is preserved individually.

\textbf{(3) $\varphi_{T'}(b) = \varphi_T(F_S(b))$ for all $b$.}
At each step $j$ reading variable $x_{i_j}$:
\begin{itemize}[nosep]
\item If $i_j \notin S$: $T_{i_j} = T'_{i_j}$ and $F_S$ does not flip $b_j$.
  Same input, same target bit, same output.
\item If $i_j \in S$: $T'_{i_j} = 1 - T_{i_j}$ and $F_S$ flips $b_j$.
  The swap of target bit and input bit cancel:
  $g_j^{(b_j)}$ under $T'$ equals $g_j^{(1-b_j)}$ under $T$, which equals
  $g_j^{(F_S(b)_j)}$ under $T$.
\end{itemize}
Hence the per-step outputs, and their product, are equal.

Combining (1)--(3):
\begin{align*}
D_{T'}(\pi) &= \Pr[\varphi_{T'}(b) = \pi \mid b \in \mathcal{I}] \\
&= \Pr[\varphi_T(F_S(b)) = \pi \mid b \in \mathcal{I}] \tag{by (3)} \\
&= \Pr[\varphi_T(b') = \pi \mid F_S^{-1}(b') \in \mathcal{I}]
\tag{substituting $b' = F_S(b)$} \\
&= \Pr[\varphi_T(b') = \pi \mid b' \in \mathcal{I}]
\tag{by (2), since $F_S = F_S^{-1}$ by (1)} \\
&= D_T(\pi). \qedhere
\end{align*}
\end{proof}

\begin{remark}
The proof uses no property of $S_5$ --- it works for branching programs of
any width over any group, provided the programs for different targets share the
same instruction template and differ only in branch assignment at the literals
whose target bit changes. Thus the argument is not specific to Barrington's
particular choice of transpositions, but it is a template-level statement rather
than a statement about arbitrary independently chosen programs.
\end{remark}

\begin{remark}
The proof extends to conjunctions with wildcards: if position $i$ is a wildcard,
the corresponding leaf outputs the same permutation $\sigma_i$ for both bit values.
Flipping $b_j$ at wildcard positions changes both branches identically, so the output
is unchanged. The bijection $F_S$ is restricted to non-wildcard differing positions.
\end{remark}

\begin{remark}[Unconditional T-independence]
\label{rem:unconditional}
The bijection $F_S$ preserves inconsistency of each variable individually
(property~(2)), and therefore also preserves the set of \emph{consistent}
paths. Consequently, $D_T = D_{T'}$ holds not only for the distribution
conditioned on inconsistent paths, but also for the unconditional distribution
over all $2^L$ paths and for the distribution over consistent paths alone.
The conditioning on inconsistency is what makes the mixing statement meaningful
--- on consistent paths the output is determined by the input, and the accepting
input is $T$ itself --- but it is not needed for T-independence, which holds for
all three laws.
\end{remark}

\begin{remark}[Template dependence]
\label{rem:grouping}
The flip bijection compares targets inside a fixed instruction template: the
same variable-read sequence, the same leaf transpositions, and the same
commutator tree. It should not be read as an invariance statement over different
templates. Changing the grouping of leaves or the orientation of internal
commutators can change the conditioned distribution, even when the resulting
program computes the same point function. The spectral results below are
therefore stated for the fixed balanced template.
\end{remark}


%% ===================================================================
\section{Spectral Contraction}
\label{sec:contraction}
%% ===================================================================

Having established that $D_T$ is $T$-independent, we now prove the harder
property: that $D_T$ is negligibly close to uniform on~$S_5$. This is where the
representation theory of $S_5$ enters, and the crux is a level-$2$-to-level-$3$
step in which a hidden simplex geometry forces contraction. Since $D_T$ is
independent of $T$ (\Cref{thm:t-indep}), it suffices to show
$\dtv(D_T, U) \leq \negl(\ell)$ for any fixed $T$; by the
Diaconis--Shahshahani bound (\Cref{thm:ds}), this reduces to bounding
$\frob{\Dhat(\rho)}^2$ for each non-trivial representation $\rho$.

\begin{lemma}[Conditioning]
\label{lem:conditioning}
The Fourier analysis below operates on the \emph{unconditional} distribution
$P_T$ in which each of the $L$ bits is drawn independently and uniformly; its
Fourier transform factorizes as a product of per-leaf projection operators
$P_\sigma$. Let $p = \Pr[b \text{ consistent}]$ and let $C_T$ denote the
distribution conditioned on consistent paths, so that
$P_T = (1-p)\,D_T + p\,C_T$. Then
\[
\dtv(D_T, U) \;\leq\; \dtv(P_T, U) + p\,\dtv(D_T, C_T)
\;\leq\; \dtv(P_T, U) + p .
\]
For the point-function program every variable is read at least once, so the
consistent paths are in bijection with $\{0,1\}^\ell$; hence
$p = 2^\ell / 2^L = 2^{-\ell(\ell-1)}$ (using $L = \ell^2$), which is negligible.
\end{lemma}

\begin{proof}
From $P_T = (1-p)\,D_T + p\,C_T$ we get $P_T - D_T = p\,(C_T - D_T)$, so
$\dtv(D_T, P_T) = p\,\dtv(D_T, C_T) \leq p$. The triangle inequality
$\dtv(D_T, U) \leq \dtv(D_T, P_T) + \dtv(P_T, U)$ then gives the bound. The count
$\abs{\{b : b \text{ consistent}\}} = 2^\ell$ holds because each of the $\ell$
variables is read at least once, so assigning a single value to each variable is a
bijection onto the consistent paths.
\end{proof}

\begin{remark}[Sign representation vanishes identically]
\label{rem:sign}
For the sign representation $(1^5)$, the leaf Fourier transform is
$\Dhat(\mathrm{sign}) = \frac{1}{2}(\mathrm{sign}(\sigma) + 1) = 0$
for any transposition~$\sigma$, since $\mathrm{sign}(\sigma) = -1$.
This propagates to all levels by multiplicativity, so the sign
representation contributes nothing to the Diaconis--Shahshahani bound.
The sum thus runs over $5$ (not $6$) non-trivial representations.
\end{remark}

\begin{proposition}[Contraction in the remaining representations]
\label{prop:other_reps}
For all five non-trivial, non-sign representations $\rho$ of $S_5$, the
operator norm $\opn{\Dhat_3(\rho)}$ is strictly less than $1$ for every
valid level-$3$ configuration. The representation $(4,1)$ is the unique
bottleneck:
\begin{center}
\begin{tabular}{lcccl}
\toprule
$\rho$ & $d_\rho$ & Leaf rank & Numerical $\max\opn{\Dhat_3(\rho)}$ & First level $< 1$ \\
\midrule
$(4,1)$     & 4 & 3 & $3.54 \times 10^{-2}$ & Level $3$ \\
$(3,2)$     & 5 & 3 & $1.93 \times 10^{-5}$ & Level $2$ \\
$(3,1,1)$   & 6 & 3 & $1.05 \times 10^{-5}$ & Level $2$ \\
$(2,2,1)$   & 5 & 2 & $< 10^{-10}$          & Level $1$ \\
$(2,1,1,1)$ & 4 & 1 & $< 10^{-10}$          & Level $1$ \\
\bottomrule
\end{tabular}
\end{center}
The representations $(2,1,1,1) = \mathrm{sign} \otimes (4,1)$ and
$(2,2,1) = \mathrm{sign} \otimes (3,2)$ have leaf projections of rank $1$
and $2$ respectively.  For rank-$1$ projections $P_\sigma = v_\sigma v_\sigma^\top$,
the level-$1$ operator $(P_\sigma P_\tau)^2 = \langle v_\sigma, v_\tau \rangle^3
\, v_\sigma v_\tau^\top$, giving
$\opn{(P_\sigma P_\tau)^2} = \abs{\langle v_\sigma, v_\tau \rangle}^3 = (1/2)^3 = 1/8$
for non-commuting pairs.  So these representations contract immediately at level~$1$. The representations $(3,2)$ and $(3,1,1)$
have rank-$3$ leaf projections like $(4,1)$, but lack the simplex
obstruction (Observation~\ref{obs:simplex}): all $405$ level-$2$ quadruples have
operator norm $< 0.085$, so they contract one level earlier.
The displayed values are independently reproduced in floating point over the
full ordered level-$3$ universe (Script~$18$). The uniform bound used in the
paper is certified separately in exact integer arithmetic for all five
representations (Script~$20$; \Cref{sec:computation}).
\end{proposition}

We focus on the bottleneck representation $(4,1)$.


\subsection{Framework: Projections and the Fourier Domain}
\label{sec:framework}

For a transposition $\sigma = (a\;b) \in S_5$, the standard representation
$\rho_{(4,1)}(\sigma)$ has eigenvalues $\{+1, +1, +1, -1\}$ (three fixed
directions and one flipped). The Fourier transform of the leaf distribution
$\{\sigma : 1/2,\; e : 1/2\}$ is
\[
P_\sigma \coloneqq \frac{1}{2}\bigl(\rho_{(4,1)}(\sigma) + I_4\bigr),
\]
which is a rank-$3$ orthogonal projection (idempotent, symmetric, trace $3$).

\begin{lemma}
\label{lem:proj_basic}
For any transposition $\sigma \in S_5$:
\begin{enumerate}[label=(\alph*),nosep]
\item $P_\sigma^2 = P_\sigma$ (idempotent),
\item $P_\sigma^\top = P_\sigma$ (symmetric),
\item $\mathrm{rank}(P_\sigma) = 3$,
\item $\frob{P_\sigma}^2 = \tr(P_\sigma) = 3$.
\end{enumerate}
\end{lemma}

The null space of $P_\sigma$ is spanned by the projection of $e_a - e_b$
onto the hyperplane $V$, where $\sigma = (a\;b)$.

Throughout this section, $\Dhat_r$ denotes the Fourier operator of a level-$r$
configuration of the \emph{unconditional} law $P_T$ (independent uniform bits at
every step); it is not the Fourier transform of the conditioned law $D_T$, which
enters only through the conditioning lemma (\Cref{lem:conditioning}).
At each level of the Barrington commutator tree, this operator is a product
of projections. For the single AND block at level~$1$ using transpositions $\sigma, \tau$
(with $\sigma^{-1} = \sigma$, $\tau^{-1} = \tau$ since they are transpositions):
\[
\Dhat_1 = P_\sigma \cdot P_\tau \cdot P_\sigma \cdot P_\tau = (P_\sigma P_\tau)^2.
\]
At level~$2$, using two level-$1$ blocks $A$ (with $\sigma_1, \tau_1$) and
$B$ (with $\sigma_2, \tau_2$):
\[
\Dhat_2 = (P_{\sigma_1} P_{\tau_1})^2 \cdot (P_{\sigma_2} P_{\tau_2})^2
\cdot (P_{\tau_1} P_{\sigma_1})^2 \cdot (P_{\tau_2} P_{\sigma_2})^2,
\]
where the inverse blocks reverse the transposition order.


\subsection{Level 0 to Level 1 Contraction}

\begin{theorem}[Universal Level-$1$ Contraction]
\label{thm:level01}
For any non-commuting pair of transpositions $(\sigma, \tau)$ in $S_5$:
\[
\frob{(P_\sigma P_\tau)^2}^2 = \frac{129}{64}.
\]
The contraction ratio is $\frac{129/64}{3} = \frac{43}{64} \approx 0.6719$,
and is \textbf{universal} --- identical for all $30$ non-commuting transposition pairs.
\end{theorem}

\begin{proof}
There are $\binom{5}{2} = 10$ transpositions in $S_5$, forming $30$
non-commuting pairs (sharing exactly one element).
For each pair, we compute $(P_\sigma P_\tau)^2$ in the orthonormal basis of
\Cref{sec:prelim} using exact symbolic arithmetic over algebraic numbers (the
basis vectors carry radicals, which are represented exactly), and verify
$\frob{(P_\sigma P_\tau)^2}^2 = 129/64$.

The universality follows from the geometry of the standard representation:
the null vectors of $P_\sigma$ and $P_\tau$ (projections of $e_a - e_b$ and
$e_c - e_d$ onto $V$) have inner product $\pm 1/2$ for all non-commuting pairs,
determined by the simplex geometry of $\{e_0, \ldots, e_4\}$ in $V$.
Since $\frob{(P_\sigma P_\tau)^2}^2$ depends only on this angle, it is
universal.
\end{proof}


\subsection{Level 1 to Level 2 Contraction}

\begin{theorem}[Level-$2$ Contraction]
\label{thm:level12}
For any valid quadruple of transpositions $(\sigma_1, \tau_1, \sigma_2, \tau_2)$
with $[\sigma_i, \tau_i] \neq e$ and $[[\sigma_1,\tau_1], [\sigma_2,\tau_2]] \neq e$:
\[
\frob{\Dhat_2}^2 \leq \frac{67586345}{67108864} \approx 1.0071.
\]
The contraction ratio relative to level $1$ is at most
$\frac{67586345/67108864}{129/64} \approx 0.4997 < 1$.
\end{theorem}

\begin{proof}
We first enumerate the reduced orientation quotient: $405$ valid quadruple
classes from the $30$ non-commuting unordered level-$1$ pairs, with the
additional constraint that the outer commutator is nontrivial. For each class,
we compute $\Dhat_2$ in exact symbolic arithmetic over algebraic numbers and
evaluate $\frob{\Dhat_2}^2$. The reduced computation yields $16$ distinct Frobenius
norms, ranging from $82431849/1073741824 \approx 0.077$ to the worst case
$67586345/67108864$.

Because the theorem is stated for ordered quadruples, we also run the ordered
level-$2$ exact cross-check (Script~$19$): $60$ ordered level-$1$ pairs give
$3240$ ordered valid level-$2$ blocks. This larger ordered universe has $17$
distinct Frobenius norms (one additional non-maximal value), but the same
maximum $67586345/67108864$.
\end{proof}


\subsection{The Simplex Obstruction}
\label{sec:simplex}

Despite the Frobenius-norm contraction at level~$2$, the operator norm
$\opn{\Dhat_2}$ equals $1$ for $270$ of the $405$ valid quadruples in the reduced census,
blocking a naive submultiplicativity argument for higher levels.
The following observation explains the structure of this obstruction.

\begin{observation}[Simplex Structure]
\label{obs:simplex}
In the reduced census, the $270$ quadruples with $\opn{\Dhat_2} = 1$
each preserve exactly one direction
in $\mathbb{R}^4$ (as the eigenvector of $\Dhat_2^\top \Dhat_2$ with eigenvalue $1$).
These directions are precisely the projections of the $5$ standard basis vectors
$e_0, \ldots, e_4 \in \mathbb{R}^5$ onto the hyperplane $V = \{\sum v_i = 0\}$:
\[
v_k = \frac{e_k - \bar{e}}{\norm{e_k - \bar{e}}}, \quad
\bar{e} = \tfrac{1}{5}(1,1,1,1,1)^\top, \quad k = 0, \ldots, 4.
\]
These $5$ directions form a \textbf{regular simplex} in $\mathbb{R}^4$, with
pairwise inner products $\abs{\langle v_i, v_j \rangle} = 1/4$ for all $i \neq j$.
Each direction is preserved by exactly $54$ quadruples (perfectly symmetric).
\end{observation}

Observation~\ref{obs:simplex} exhausts this reduced census. The following lemma supplies the
structural implication that the level-$3$ argument actually uses: preserving
$v_k$ forces every leaf of the block to fix the point $k$.

\begin{lemma}[Stabilizer characterization of preserved directions]
\label{lem:stab-char}
Let $B$ be a valid level-$2$ block with leaf transpositions
$\sigma_1, \tau_1, \sigma_2, \tau_2$, and let $\Dhat_2(B)$ be its level-$2$
operator in the standard representation. For each $k \in \{0,\dots,4\}$:
\[
\norm{\Dhat_2(B)\, v_k} = \norm{v_k}
\qquad\Longleftrightarrow\qquad
\sigma_i(k) = \tau_i(k) = k \ \ (i = 1,2),
\]
and in that case $\Dhat_2(B)\, v_k = v_k$. Equivalently: a level-$2$ block
preserves the direction $v_k$ if and only if all four of its leaf
transpositions lie in the point stabilizer $\mathrm{Stab}(k) \cong S_4$.
\end{lemma}

\begin{proof}
Write $\Dhat_2(B) = P_{t_1} P_{t_2} \cdots P_{t_r}$ as the product of the
$r = 16$ leaf projections appearing in its expansion (each factor is
$P_t = \frac{1}{2}(\rho_{(4,1)}(t) + I)$ for one of the four leaf
transpositions~$t$; the inverse sub-blocks contribute the same projections in
reversed order, since each transposition is an involution). By
\Cref{lem:proj_basic} every $P_t$ is an orthogonal projection, hence
$\norm{P_t w} \leq \norm{w}$ for all $w$.

$(\Rightarrow)$ Suppose $\norm{\Dhat_2(B)v_k} = \norm{v_k}$. Setting
$w_0 = v_k$ and $w_j = P_{t_{r-j+1}} w_{j-1}$, the chain
\[
\norm{v_k} = \norm{w_r} \leq \norm{w_{r-1}} \leq \cdots \leq \norm{w_0}
= \norm{v_k}
\]
collapses to equality at every step. For an orthogonal projection,
$\norm{P w} = \norm{w}$ forces $Pw = w$ (indeed
$\norm{w}^2 = \norm{Pw}^2 + \norm{(I-P)w}^2$, so $\norm{(I-P)w} = 0$).
Applying this successively gives $P_{t}\,v_k = v_k$ for every leaf
transposition $t$ occurring in the product.

Now $P_t v_k = v_k$ means $\rho_{(4,1)}(t)\,v_k = v_k$. Since $v_k$ is the
projection of the basis vector $e_k$ onto the sum-zero hyperplane $V$ and
$\rho_{(4,1)}$ is the restriction of the permutation action to $V$, we have
$\rho_{(4,1)}(t)\,v_k = v_{t(k)}$. The five vectors $v_0,\dots,v_4$ are
pairwise distinct (they satisfy $\langle v_i, v_j\rangle = -1/4 \neq 1$ for
$i \neq j$), so $v_{t(k)} = v_k$ forces $t(k) = k$. Hence every leaf
transposition of $B$ fixes $k$, i.e.\ all four lie in $\mathrm{Stab}(k)$.

$(\Leftarrow)$ If all four leaf transpositions fix $k$ then
$\rho_{(4,1)}(t) v_k = v_{t(k)} = v_k$ for each of them, so $P_t v_k = v_k$ and
therefore $\Dhat_2(B) v_k = v_k$, which has the same norm as $v_k$.
\end{proof}

The regularity of this structure --- five directions, uniform distribution,
constant pairwise angle --- might suggest a general phenomenon. The following
theorem shows the opposite: the simplex structure is specific to $S_5$, and
for $S_n$ with $n \geq 6$ a dimensional obstruction prevents the same
level-$2$ direction-mixing argument from proving a uniform level-$3$ gap.

\begin{theorem}[Dimensional obstruction for $S_n$, $n \geq 6$]
\label{thm:dim-obstruction}
For $S_n$ with $n \geq 5$, the level-$1$ Fourier transform
$\Dhat_1 = (P_\sigma P_\tau)^2$ in the standard representation $(n-1,1)$
has singular values $\underbrace{1, \ldots, 1}_{n-3},\; 1/8,\; 0$.
In particular, $\Dhat_1$ preserves a subspace of dimension $n-3$ isometrically.

For $n \geq 6$, the minimum intersection of two such preserved subspaces
in $\mathbb{R}^{n-1}$ is $2(n-3)-(n-1) = n-5 \geq 1$, so every valid level-$2$
quadruple has $\opn{\Dhat_2} = 1$ and the simplex-breaking mechanism of
\Cref{lem:mixing} does not extend by the same argument. At level~$3$, the
maximum operator norm is~$1$ for $S_6$ and $S_7$, whereas every configuration
has norm~$1$ for $n \geq 10$. Thus $S_5$ is the only case with uniform
level-$3$ contraction among the cases settled here; the present analysis does
not decide $n \in \{8,9\}$.
\end{theorem}

\begin{proof}
The projections $P_\sigma = I_V - v_\sigma v_\sigma^\top$ and
$P_\tau = I_V - v_\tau v_\tau^\top$ (where $v_\sigma, v_\tau$ are the unit
null vectors of the transpositions in~$V$) share the property that
$\langle v_\sigma, v_\tau \rangle = -1/2$, after choosing their signs, for all non-commuting pairs
(independent of~$n$). Decomposing $V = W \oplus W^\perp$ where
$W = \mathrm{span}(v_\sigma, v_\tau)$ has dimension~$2$: on $W^\perp$
(dimension $n-3$), both projections act as the identity, so
$\Dhat_1 = (P_\sigma P_\tau)^2 = I$ on $W^\perp$, contributing $n-3$
singular values equal to~$1$. On~$W$, the restriction is a $2 \times 2$
block independent of~$n$; a direct computation gives
$\Dhat_1^\top \Dhat_1\big|_W$ with eigenvalues $1/64$ and~$0$
(trace $= 4/256$, determinant $= 0$), so the singular values on $W$
are $1/8$ and~$0$.

For $n \geq 6$: two subspaces of dimension $n-3$ in $\mathbb{R}^{n-1}$
intersect in dimension $\geq 2(n-3) - (n-1) = n-5 \geq 1$. Any unit vector
in this intersection is preserved isometrically by both level-$1$ sub-blocks
of a level-$2$ commutator, so $\opn{\Dhat_2} = 1$ for every valid quadruple.

The same counting one step up controls level~$3$. Each level-$2$ operator with
$\opn{\Dhat_2} = 1$ fixes a subspace of dimension $\geq n-5$, and two such
subspaces in $\mathbb{R}^{n-1}$ share a common fixed vector once
$2(n-5)-(n-1) = n-9 \geq 1$, i.e.\ for $n \geq 10$; then $\opn{\Dhat_3} = 1$ as
well. The value $n = 9$ is exactly where disjoint placement becomes possible
again, so the dimensional argument alone does not settle $n \in \{8,9\}$. For the
two immediate neighbours, the depth-$3$ seed given after
\Cref{def:validtemplate}, viewed in $S_6$ or $S_7$, remains valid and fixes
the last point. Every leaf projection therefore fixes the nonzero vector
$(1,\ldots,1,-(n-1))\in V$, so its level-$3$ operator has norm~$1$.
Script~$14$ also checks the maximum numerically over the reduced census of
$1{,}242{,}915$ configurations in $S_6$ and $10{,}282{,}335$ in $S_7$.
This is a statement about the maximum, not every configuration: the script
includes exact Gram-form certificates for strictly contractive examples in
both groups. Thus uniform level-$3$ contraction fails in these two cases.
\end{proof}

\begin{remark}
The one-dimensional simplex mechanism of Observation~\ref{obs:simplex} does not
generalize cleanly to $S_n$ for $n \geq 6$: exhaustive computation confirms
that all $1620$ quadruples in the reduced census for $S_6$ and all $4725$ for $S_7$ have
operator norm $1$ (versus $270/405 = 66.7\%$ for $S_5$).  In these larger
groups the preserved isometric subspaces have forced positive-dimensional
intersections rather than the disjoint direction placement that drives the
$S_5$ argument. \Cref{thm:dim-obstruction} explains this: $S_5$ is the unique
case where the preserved subspace dimension ($n-3 = 2$) is small enough
relative to the ambient dimension ($n-1 = 4$) to permit disjoint placement.
The theorem is verified by two independent methods, including exact symbolic
arithmetic over all adjacent pairs in $S_5, S_6, S_7$ and exhaustive
enumeration of the reduced level-$2$ censuses (Scripts~$14$ and~$15$).
\end{remark}

An immediate consequence is that the level-$3$ spectral contraction mechanism
presented here does not directly extend to the width $> 5$ programs that appear
in optimized obfuscation constructions avoiding Barrington's $4^d$
blowup~\cite{AGIS14}. Establishing unconditional mixing for such programs
requires fundamentally different techniques, which we discuss in
\Cref{sec:extensions}.


\subsection{Level 2 to Level 3 Contraction}

\paragraph{Why switch to the operator norm.}
The Frobenius norm tracks total ``energy'' of the Fourier operator and is
computed exactly at the first two levels (Theorems~\ref{thm:level01}--\ref{thm:level12});
the induction from level~$3$ onward uses the operator norm.
However, the worst-case $\frob{\Dhat_2}^2 \approx 1.007 > 1$, so the
absolute Frobenius norm has not yet dropped below~$1$, and an inductive
argument based on Frobenius norms alone does not yield $\frob{\Dhat_d} \to 0$.
The operator norm resolves this: once $\opn{\Dhat_3} < 1$,
submultiplicativity raises the uniform level-$d$ bound to its fourth power,
which iterates to doubly-exponential decay in the depth.  The cost is that we need
operator-norm contraction at level~$3$ for \emph{every} configuration,
not just worst-case Frobenius contraction.

\begin{lemma}[Direction Mixing]
\label{lem:mixing}
For every valid level-$3$ configuration $(\mathrm{quad}_A, \mathrm{quad}_B)$
in which both sub-blocks have $\opn{\Dhat_2} = 1$, the two sub-blocks
preserve \textbf{different} directions. In the reduced census, there are
$29160$ such pairs, none sharing the same preserved direction.
\end{lemma}

\begin{proof}
Suppose for contradiction that both $\mathrm{quad}_A$ and $\mathrm{quad}_B$
preserve the same direction $v_k$ (the projection of $e_k$ onto $V$).
By \Cref{lem:stab-char}, a quadruple preserving $v_k$ consists of transpositions
that fix element $k$, so all four transpositions in each sub-block lie in
the stabilizer $\mathrm{Stab}(k) \cong S_4$.

The derived series of $S_4$ is
$S_4 \rhd A_4 \rhd V_4 \rhd \{e\}$,
where $V_4 = \{e, (01)(23), (02)(13), (03)(12)\}$ is the Klein four-group
(abelian).  In the Barrington commutator tree:
\begin{itemize}[nosep]
\item Level-$1$ targets $[\sigma_i, \tau_i]$ are commutators of transpositions
  in $S_4$, hence lie in $[S_4, S_4] = A_4$.
\item Level-$2$ targets $[[\sigma_1,\tau_1],[\sigma_2,\tau_2]]$ are
  commutators of elements of $A_4$, hence lie in $[A_4, A_4] = V_4$.
\end{itemize}
Since $V_4$ is abelian, the outer commutator
$[\alpha_A, \alpha_B] = e$ for any $\alpha_A, \alpha_B \in V_4$.
This contradicts the requirement $[\alpha_A, \alpha_B] \neq e$ for a
valid level-$3$ configuration.

The result is also confirmed by exhaustive enumeration over the reduced
configuration set (Script~$11$).
\end{proof}

\begin{theorem}[Level-$3$ Contraction]
\label{thm:level23}
For every valid level-$3$ configuration:
\[
\opn{\Dhat_3} \leq \lambda \coloneqq 0.0355.
\]
In particular, $\opn{\Dhat_3} < 1$ for all configurations in the fixed
transposition template.
\end{theorem}

\begin{proof}
For each valid level-$3$ pair $(\mathrm{quad}_A, \mathrm{quad}_B)$, we compute
\[
\Dhat_3 = \Dhat_{2,A} \cdot \Dhat_{2,B} \cdot \Dhat_{2,A}^\top \cdot \Dhat_{2,B}^\top
\]
over the full ordered universe: all five non-trivial, non-sign irreducible
representations over the $3{,}240$ ordered level-$2$ blocks, i.e.\ $9{,}447{,}840$
oriented pairs. Script~$20$ reduces this universe, with exact orbit accounting,
to $39{,}591$ representatives under simultaneous $S_5$-conjugation and exchange
of the two blocks.

The five relevant irreducibles are carried by four integral carrier modules
derived from permutation actions (with $M_9^{\mathrm{s}}$ obtained by a sign
twist), so that no irrational basis is ever needed. Let $M_{41}$ be the sum-zero
submodule of the natural action on the five points, which realises $(4,1)$; let
$M_9$ be the sum-zero submodule of the action on the ten $2$-subsets, whose
decomposition is $(4,1)\oplus(3,2)$; let $M_{9}^{\mathrm{s}}$ be the sign twist
of $M_9$, carrying $(2,1,1,1)\oplus(2,2,1)$; and let
$M_6 = \textstyle\bigwedge^2 M_{41}$, which realises $(3,1,1)$. The invariant
Gram forms are $G = I + J$ on $M_{41}$, $M_9$, and $M_9^{\mathrm{s}}$ (with $J$
the all-ones matrix on the ambient coordinates) and $\bigwedge^2(I+J)$ on $M_6$.
Each of these decompositions is multiplicity-free, so the four carriers cover
the five relevant irreducibles and an operator-norm bound on each carrier is
exactly a bound on the irreducibles it contains. With $\Dhat_3 = N/2^{64}$ for
an integer matrix $N$, the script verifies in exact integer arithmetic that
\[
  71^2\,2^{128}G-2000^2N^\top G N \succ 0
\]
for every representative, using fraction-free Bareiss leading principal
minors. Hence $\opn{\Dhat_3}\le71/2000=0.0355$ exactly. A second all-irrep pass
certifies the tighter global upper bound $0.0354002321$. For the bottleneck
configuration, which lives on $M_{41}$, the script forms $S = N^\top G N$ and
the pencil $p(y) = \det(yG - S)$; a Sturm chain isolates the largest root
$y_{\max}$ of $p$ in exact rational arithmetic, and
$\sigma_{\max}^2 = y_{\max}/2^{128}$, which gives
\[
  0.035400232083798851\ldots
  \le \max\opn{\Dhat_3} \le 0.0354002321.
\]
The full certificate format---carriers, Gram forms, orbit accounting, and the
integer pencil---is documented with the artifact
(\texttt{CERTIFICATE\_LEVEL3.md}). Scripts~$17$ and~$18$ independently
reproduce the same bottleneck numerically.

The contraction occurs because:
\begin{itemize}[nosep]
\item If at least one sub-block has $\opn{\Dhat_2} < 1$, submultiplicativity gives
  $\opn{\Dhat_3} \leq \opn{\Dhat_{2,A}}^2 \cdot \opn{\Dhat_{2,B}}^2 < 1$.
\item If both sub-blocks have $\opn{\Dhat_2} = 1$ but preserve different directions
  $v_A \neq v_B$ (guaranteed by \Cref{lem:mixing}), then $\Dhat_{2,A}$ maps $v_B$
  to a vector with component $\abs{\langle v_A, v_B \rangle} = 1/4$ along $v_A$,
  causing the product to contract. \qedhere
\end{itemize}
\end{proof}


\subsection{Induction to All Levels}

\begin{theorem}[Exponential Contraction]
\label{thm:induction}
For every fixed valid balanced transposition-template program for point functions
$\CC_T$ over $\{0,1\}^\ell$ with $\ell = 2^d$, $d \geq 3$:
\[
\opn{\Dhat_d\bigl((4,1)\bigr)} \leq \lambda^{4^{d-3}},
\]
where $\lambda = 0.0355$.
\end{theorem}

\begin{proof}
At depth $d+1$, the commutator structure gives
$\Dhat_{d+1} = \Dhat_A \cdot \Dhat_B \cdot \Dhat_A^\top \cdot \Dhat_B^\top$.
The sub-blocks are valid whenever the root configuration is valid: if a
sub-block target were $e$, its ancestor commutator would be $[e,\cdot]=e$,
forcing the root target to be trivial. Thus the level-$d$ bound applies to
both children.
By submultiplicativity of the operator norm and $\opn{M^\top} = \opn{M}$:
\[
\begin{aligned}
\opn{\Dhat_{d+1}}
&\leq \opn{\Dhat_A} \cdot \opn{\Dhat_B} \cdot
        \opn{\Dhat_A^\top} \cdot \opn{\Dhat_B^\top} \\
&= \opn{\Dhat_A}^2\,\opn{\Dhat_B}^2 \\
&\leq \bigl(\lambda^{4^{d-3}}\bigr)^4
 = \lambda^{4^{(d+1)-3}}.
\end{aligned}
\]
The last inequality applies the uniform level-$d$ induction hypothesis
separately to the potentially different children $A$ and~$B$.
The base case $\opn{\Dhat_3} \leq \lambda$ is \Cref{thm:level23}.
Iterating: $\opn{\Dhat_d} \leq \lambda^{4^{d-3}}$.
\end{proof}


\subsection{\texorpdfstring{Why $S_5$ is Structurally Unique}{Why S5 is Structurally Unique}}
\label{sec:s5-unique}

\Cref{thm:level23,thm:dim-obstruction} together provide a precise structural
explanation of why $S_5$ is, among the symmetric groups we settle, the unique
one where uniform level-$3$ spectral contraction succeeds in the standard representation.
Three properties conspire:

\begin{enumerate}[nosep]
\item \textbf{$S_5$ is non-solvable.} This makes it possible to choose
  non-commuting targets and conjugates so that the Barrington commutator
  construction does not collapse to the identity. It does not mean that
  arbitrary non-identity elements have nontrivial commutator.
\item \textbf{$\mathrm{Stab}(k) \cong S_4$ is solvable with derived length
  exactly~$3$.} The derived series
  $S_4 \rhd A_4 \rhd V_4 \rhd \{e\}$ terminates at depth~$3$, which is
  precisely the level of the commutator tree where contraction occurs.
  This forces the direction-mixing mechanism of \Cref{lem:mixing}.
\item \textbf{The preserved subspace dimension is at the critical half
  threshold.} $\Dhat_1$ preserves $n-3 = 2$ of $n-1 = 4$
  dimensions, exactly half of the ambient space, allowing two preserved
  subspaces to be disjoint. For $n \geq 6$, the ratio
  $(n-3)/(n-1) > 1/2$ makes disjoint placement impossible
  (\Cref{thm:dim-obstruction}).
\end{enumerate}

These properties are not independent coincidences. The first two express a
single structural condition on the point stabilizer: a level-$2$ block preserves
a direction $v_k$ exactly when its transpositions fix $k$---that is, lie in
$\mathrm{Stab}(k) \cong S_{n-1}$ (\Cref{lem:stab-char})---and the direction-mixing
mechanism that breaks $v_k$ at level~$3$ is driven by the \emph{solvability} of
$S_{n-1}$: for $S_5$, $\mathrm{Stab}(k) \cong S_4$ is solvable and its derived
length ($3$) is exactly the level at which contraction occurs. $S_5$ is the
largest symmetric group whose point stabilizer is solvable, and the settled cases
fit this picture---$S_6, S_7$ (non-solvable stabilizer) admit norm-$1$
configurations, and $n \geq 10$ is ruled out dimensionally. The present
analysis does not decide $n \in \{8,9\}$. This
points past the symmetric groups: it suggests that commutator-tree mixing in a
finite group may be governed not by the group's size but by the solvability of
the subgroups its leaves generate---a possible solvability barrier. Any extension
of the present framework---to wider symmetric groups, or to
any other finite group---must therefore confront this solvability structure
rather than the size of the group, and must in each case specify the leaf class,
the representation analysed, the fixed spaces of the resulting operators, and
the subgroups the leaves generate.


%% ===================================================================
\section{Main Result}
\label{sec:main}
%% ===================================================================

We now assemble the pieces---exact $T$-independence (\Cref{thm:t-indep}), the
conditioning lemma (\Cref{lem:conditioning}), and the level-$3$ contraction with
its inductive amplification (\Cref{thm:level23,thm:induction})---into the two
quantitative guarantees.

\begin{theorem}[Main Theorem]
\label{thm:main}
For every fixed valid balanced transposition-template program for point functions
$\CC_T$ over $\{0,1\}^\ell$ with $\ell = 2^d$, $d \geq 3$:
\begin{enumerate}[label=(\alph*),nosep]
\item $D_T = D_{T'}$ for all $T, T' \in \{0,1\}^\ell$.
\item $\dtv(D_T, \mathrm{Uniform}(S_5)) \leq 6 \cdot \lambda^{4^{d-3}} + 2^{-\ell(\ell-1)}$,
  where $\lambda = 0.0355$ (the additive term is the conditioning loss of
  \Cref{lem:conditioning}).
\item In particular, $\dtv(D_T, U) = \negl(\ell)$.
\end{enumerate}
\end{theorem}

\begin{proof}
Part (a) is \Cref{thm:t-indep}. For part (b), the Fourier transform $\Dhat$
computed below is that of the unconditional distribution $P_T$, so by
\Cref{thm:ds}:
\[
\dtv(P_T, U)^2 \leq \frac{1}{4} \sum_{\rho \neq \mathrm{triv}} d_\rho \cdot
\frob{\Dhat(\rho)}^2.
\]
By \Cref{rem:sign}, the sign representation contributes zero.
For each remaining representation $\rho$ of dimension $d_\rho$,
\Cref{prop:other_reps} gives a base case
$\opn{\Dhat_3(\rho)} \leq \lambda_\rho < 1$ for all valid level-$3$
configurations. The submultiplicativity argument of \Cref{thm:induction}
applies to each $\rho$ independently (it uses only $\opn{M^\top} = \opn{M}$
and submultiplicativity, which hold in any representation), yielding
$\opn{\Dhat_d(\rho)} \leq \lambda_\rho^{4^{d-3}}$.
Since $\frob{M}^2 \leq d_\rho \cdot \opn{M}^2$:
\[
\sum_{\rho} d_\rho \cdot \frob{\Dhat(\rho)}^2
\leq \sum_{\rho} d_\rho^2 \cdot \lambda_\rho^{2 \cdot 4^{d-3}}
\leq \Bigl(\sum_{\rho} d_\rho^2\Bigr) \cdot \lambda^{2 \cdot 4^{d-3}},
\]
where the exact global certificate gives
$\max_\rho\lambda_\rho\le0.0354002321<0.0355$ and the standard representation
contains the certified lower endpoint $0.035400232083798851\ldots$
(\Cref{prop:other_reps} and \Cref{thm:level23}). We use the conservative
$\lambda=0.0355$. Since
$\sum_\rho d_\rho^2 = 4^2 + 5^2 + 6^2 + 5^2 + 4^2 = 118$:
\[
\dtv(P_T, U)^2 \leq \frac{118}{4} \cdot \lambda^{2 \cdot 4^{d-3}},
\qquad
\dtv(P_T, U) \leq \sqrt{\tfrac{118}{4}} \cdot \lambda^{4^{d-3}}
< 6 \cdot \lambda^{4^{d-3}}.
\]
Finally, by \Cref{lem:conditioning},
\[
\dtv(D_T, U) \leq \dtv(P_T, U) + 2^{-\ell(\ell-1)}
< 6 \cdot \lambda^{4^{d-3}} + 2^{-\ell(\ell-1)}.
\]
For part (c), since $\lambda < 1$ and $4^{d-3}$ grows exponentially in
$d = \log_2 \ell$ (so $\lambda^{4^{d-3}}$ decays doubly-exponentially in $d$,
i.e. like $\exp(-\Theta(\ell^2))$ in $\ell$), the bound is negligible in $\ell$.
\end{proof}

\subsection{The fresh-sample terminal-output experiment}
\label{subsec:terminal}

The Main Theorem controls one exactly specified experiment. We record its
$q$-fold form, since that is the form a reduction would consume, and then
delimit precisely what it does and does not cover.

\begin{definition}[Fresh-sample terminal-output transcript]
\label{def:terminal}
Fix $\ell = 2^d$ with $d \geq 3$ and a target $T \in \{0,1\}^\ell$. The
\emph{$q$-sample terminal-output transcript} of $\BP_T$ is
$(\pi_1, \dots, \pi_q)$, where each $\pi_i \in S_5$ is the terminal output
permutation of $\BP_T$ evaluated on a \emph{fresh, independent, uniformly
random inconsistent path}; the $\pi_i$ are thus i.i.d.\ with law $D_T$, and we
write $D_T^{\otimes q}$ for the law of the transcript.
\end{definition}

\begin{corollary}[Terminal-output indistinguishability]
\label{cor:terminal}
Let $\varepsilon(\ell) = 6\lambda^{4^{d-3}} + 2^{-\ell(\ell-1)}$ be the bound of
\Cref{thm:main}(b), with $\lambda = 0.0355$. For every integer
$q \in \mathbb{Z}_{\geq 1}$ and all $T, T' \in \{0,1\}^\ell$:
\begin{enumerate}[label=(\alph*),nosep]
\item $D_T^{\otimes q} = D_{T'}^{\otimes q}$ (exactly, not merely
  statistically);
\item $\dtv\bigl(D_T^{\otimes q},\, \mathrm{Uniform}(S_5)^{\otimes q}\bigr)
  \leq \min\{1,\; q\,\varepsilon(\ell)\}$;
\item for every target-independent randomised post-processing $K$ --- a Markov
  kernel from $S_5^q$ to an arbitrary finite space, fixed without knowledge of
  $T$ --- both conclusions persist:
  $K D_T^{\otimes q} = K D_{T'}^{\otimes q}$ and
  $\dtv\bigl(K D_T^{\otimes q},\, K\,\mathrm{Uniform}(S_5)^{\otimes q}\bigr)
  \leq \min\{1,\; q\,\varepsilon(\ell)\}$.
\end{enumerate}
\end{corollary}

\begin{proof}
(a) By \Cref{thm:main}(a) the coordinate laws coincide, $D_T = D_{T'}$; the
coordinates are independent, hence the product laws coincide.

(b) Let $H_j$ be the hybrid law whose first $j$ coordinates are uniform on $S_5$
and whose remaining $q-j$ coordinates are i.i.d.\ $D_T$, so that
$H_0 = D_T^{\otimes q}$ and $H_q = \mathrm{Uniform}(S_5)^{\otimes q}$.
Consecutive hybrids differ in one coordinate only, and total variation between
product measures differing in a single independent coordinate equals the total
variation of that coordinate, so $\dtv(H_j, H_{j+1}) =
\dtv(D_T, \mathrm{Uniform}(S_5)) \leq \varepsilon(\ell)$ by \Cref{thm:main}(b).
Telescoping with the triangle inequality over $j = 0, \dots, q-1$ gives
$\dtv(D_T^{\otimes q}, \mathrm{Uniform}(S_5)^{\otimes q}) \leq q\,\varepsilon(\ell)$;
the bound $1$ is trivial.

(c) Total variation contracts under Markov kernels (data processing):
$\dtv(K\mu, K\nu) \leq \dtv(\mu, \nu)$ for every kernel $K$ and all laws
$\mu, \nu$. Since $K$ is fixed independently of $T$, part (a) gives
$K D_T^{\otimes q} = K D_{T'}^{\otimes q}$, and data processing applied to
part (b) gives the stated bound.
\end{proof}

\begin{remark}[Scope of \Cref{cor:terminal}]
\label{rem:terminal-limits}
\Cref{cor:terminal} concerns the terminal-output marginal under \emph{fresh,
independent, uniformly random} inconsistent paths, together with
target-independent post-processing of that transcript. Its hypotheses distinguish
this experiment from three other kinds of observation:
\begin{itemize}[nosep,leftmargin=1.2em]
\item fixed, adversarially chosen, or adaptive paths, including choices based
  on earlier outputs;
\item correlated evaluations, such as evaluations sharing path bits;
\item intermediate products or an encoded adversarial view, including Kilian
  randomisation, encoded matrices, and graded encodings.
\end{itemize}
Applying the result to any of these settings requires an additional argument
relating that observation to the stated experiment. The construction-level
example in \Cref{sec:intro} describes this requirement for obfuscation.
\end{remark}

\begin{remark}[Non-power-of-$2$ input lengths]
\label{rem:padding}
The theorem is stated only for the complete balanced tree with $\ell = 2^d$.
Padding by leaves that output the identity in both branches is not innocuous in
a commutator tree: a subcommutator of the form $[g,e]$ collapses to the identity
and can change the target of the parent subtree. Extending the result to
arbitrary $\ell$ therefore requires a separate template and a separate count of
consistent and inconsistent paths.
\end{remark}

For concrete values:
\begin{center}
\begin{tabular}{cccc}
\toprule
$\ell$ & Depth $d$ & $\opn{\Dhat_{(4,1)}}$ bound & $\dtv$ bound \\
\midrule
$8$   & $3$ & $3.55 \times 10^{-2}$ & $2.13 \times 10^{-1}$ \\
$16$  & $4$ & $1.59 \times 10^{-6}$ & $9.53 \times 10^{-6}$ \\
$32$  & $5$ & $6.36 \times 10^{-24}$ & $3.82 \times 10^{-23}$ \\
$64$  & $6$ & $1.64 \times 10^{-93}$ & $9.83 \times 10^{-93}$ \\
\bottomrule
\end{tabular}
\end{center}


%% ===================================================================
\section{Extensions and Open Problems}
\label{sec:extensions}
%% ===================================================================

\subsection{Conjunctions with Wildcards}

A conjunction $\CC_{T,W}$ with target $T \in \{0,1\}^\ell$ and wildcard set
$W \subseteq [\ell]$ outputs $1$ iff $x_i = T_i$ for all $i \notin W$.
In the branching program, wildcard positions output $\sigma_i$ for both
bit values.  Obfuscation of conjunctions has been studied
extensively~\cite{BVWW16,BKMPRS18,BLMZ19}, which is what motivates looking at
this class here. For conjunctions we prove T-independence of the terminal
marginal and report computational evidence for mixing.

\Cref{thm:t-indep} extends directly to conjunctions:
$D_{T,W} = D_{T',W}$ for all $T, T'$ with the same wildcard set $W$.
The spectral contraction depends on the wildcard density $|W|/\ell$.
Computational evidence shows strong contraction for $|W|/\ell \leq 1/2$,
weakening contraction for $|W|/\ell = 3/4$, and no contraction for $|W|/\ell = 1$
(the trivial function). The precise threshold is an open problem.


\subsection{Beyond Point Functions}

For general branching programs computing the same Boolean function via
different constructions, ``strong universal mixing'' ($D_{C_0} = D_{C_1}$)
is \textbf{false}: programs with different transposition assignments produce
different distributions. However, ``weak universal mixing'' (both distributions
converge to uniform independently) appears to hold: computational evidence shows
equivalent circuits with different pair assignments achieving similar $\dtv$
values. Proving this would extend the algebraic terminal-output property to a
larger class of programs.

Computational experiments reveal a stronger pattern.  For $\ell = 4$, two
programs computing $\mathrm{AND}_4$ with different transposition pairs
yield Fourier norms $\frob{\Dhat(4,1)}^2 = 0.147$ and $0.293$
respectively --- clearly different distributions --- yet \emph{identical}
total variation distances $\dtv = 0.309$ to uniform.  A sampled OR-of-AND
experiment shows approximately equal distances for two structurally different
implementations ($\dtv=0.1904$ versus $0.1923$). If this ``equidistance''
phenomenon is general,
it would mean that $\dtv(D_C, U)$ depends only on the Boolean function
computed, not on the specific Barrington construction --- a property
strictly stronger than weak universal mixing.


\subsection{Generalization to Other Groups}

The dimensional obstruction in \Cref{thm:dim-obstruction} prevents a direct
extension of the $S_5$ direction-mixing proof to $S_n$ with $n \geq 6$.
This motivates studying other representations, deeper levels of the
commutator tree, and other groups. Finite matrix groups such as
$\mathrm{GL}_n(\mathbb{F}_q)$ are natural
mathematical candidates for such a move. The obstruction of
\Cref{thm:dim-obstruction} is specific to the transposition leaf class and to
the standard representation of $S_n$; it therefore does not decide the matrix
case, which would first require fixing a leaf class and a representation before
the analogous fixed-space count could be stated.

Our proof follows a three-step program that is, in principle, portable:
(1)~define the output distribution on inconsistent inputs,
(2)~identify the bottleneck irreducible representation,
(3)~discover geometric structure in the spectral obstructions.
For $\mathrm{GL}_2(\mathbb{F}_q)$ with $q$ a small prime, the irreducible
representations are completely classified (principal series, cuspidal,
Steinberg, and one-dimensional characters), and the analogue of step~(2) is
computationally feasible.  Determining whether the level-$2$ obstructions
for matrix branching programs over $\mathrm{GL}_2(\mathbb{F}_q)$ exhibit
geometric regularity analogous to the simplex in Observation~\ref{obs:simplex} would
be a concrete first test of whether spectral contraction extends beyond
symmetric groups.

Establishing contraction in every nontrivial irreducible would yield an
unconditional mixing theorem for such a template over a matrix group.
A cryptographic application would then require the construction-specific
connection described in \Cref{sec:intro}.


\subsection{Deeper Levels and the Solvability Barrier}

Two independent obstructions make $S_5$ special in the standard
representation, and extending beyond $S_5$ requires confronting both.

The first, discussed in \Cref{sec:s5-unique}, is dimensional
(\Cref{thm:dim-obstruction}): for $n \geq 6$, the preserved subspace of
$\Dhat_1$ has dimension $n-3$ in $\mathbb{R}^{n-1}$, forcing every
level-$2$ quadruple to have $\opn{\Dhat_2} = 1$, which blocks level-$3$
operator-norm contraction from following the same direction-mixing proof used
for $S_5$.

The second is algebraic. The proof of \Cref{lem:mixing} relies on the
solvability of~$S_4$: the derived series $S_4 \rhd A_4 \rhd V_4 \rhd \{e\}$
forces level-$2$ targets into $V_4$ (abelian), killing the outer commutator.
For $S_n$ with $n \geq 6$, the relevant stabilizer is $S_{n-1}$, which is
\emph{non-solvable}, so even if the dimensional obstruction could be
sidestepped, the direction-mixing argument would require a fundamentally
different mechanism.

The analysis presented here leaves the following questions undecided.
For $S_6$, $S_7$, and $n \geq 10$, where uniform contraction in the standard
representation fails at level~$3$, one may ask whether a uniform bound
$\opn{\Dhat_d} < 1$ holds at some level $d \geq 4$. For $S_8$ and $S_9$,
the corresponding question starts at level~$3$. A separate path, for
every $n \geq 6$, is to study \emph{non-standard representations}: whether some
irreducible representation other than $(n-1,1)$ contracts at level~$3$. Note that a
Diaconis--Shahshahani mixing bound requires controlling \emph{all} nontrivial
irreducible representations simultaneously, so contraction in a single
non-standard representation alone does not yield mixing; the bottleneck may
differ from $(n-1,1)$. These constitute the most natural follow-up directions.


%% ===================================================================
\section{Computational Verification}
\label{sec:computation}
%% ===================================================================

The proofs of \Cref{thm:level01,thm:level12} use exact symbolic algebraic
arithmetic (implemented in \texttt{sympy}) with a verified orthonormal basis for
the standard representation: that basis carries the radicals $\sqrt{2}$,
$\sqrt{6}$, $\sqrt{12}$, $\sqrt{20}$, which \texttt{sympy} represents exactly, so
the intermediate quantities are exact algebraic numbers (the reported Frobenius
norms are rational). \Cref{thm:level23} uses an exact integer Gram-form
certificate over the full ordered universe of level-$3$ configurations---
$9{,}447{,}840$ oriented pairs across all five non-trivial, non-sign
representations, partitioned into $39{,}591$ symmetry orbits with exact accounting
(Script~$20$). Scripts~$17$ and~$18$ provide independent full-universe
\texttt{numpy} cross-checks.

The reproducibility archive is available at
\url{https://github.com/ernest-zermeno/uniform-mixing-barrington-branching-programs}.
It contains all $20$ verification
scripts, the exact level-$3$ certificate, analytic proof notes, dependencies,
and instructions for reproducing the checks.

\begin{center}
\begin{tabular}{clc}
\toprule
Script & Content & Method \\
\midrule
01 & Ambivalence of $S_5$ (similarity blindness) & Exhaustive \\
02 & Collision analysis, $\ell = 4$ & Exhaustive \\
03 & Collision scaling, $\ell = 4 \to 8$ & Sampled \\
04 & Fourier spectral analysis & Exact + sampled \\
05 & Barrington construction + Kilian obfuscation & Exhaustive \\
06 & Contraction operator (unconditional distribution) & Analytic \\
07 & T-independence proof verification & Exhaustive \\
08 & Conjunctions with wildcards & Exhaustive \\
09 & Universal mixing investigation & Exact + sampled \\
10 & Complete spectral contraction proof & Exact + numerical \\
11 & Eigenvector analysis (simplex discovery) & Numerical \\
12 & Reduced level-$3$ census, standard representation & Numerical \\
13 & Reduced level-$3$ census, all representations & Numerical \\
14 & Obstruction checks ($S_5, S_6, S_7$) & Numerical + exact \\
15 & Exact sympy proof of singular values of $\Dhat_1$ & Exact \\
16 & Conditioning lemma, $\ell = 4$ & Exact \\
17 & Ordered-orientation level-$3$ cross-check & Numerical \\
18 & All-irreps ordered-orientation cross-check & Numerical \\
19 & Ordered level-$2$ exact cross-check & Exact \\
20 & Full-universe level-$3$ Gram certificate & Exact integer \\
\bottomrule
\end{tabular}
\end{center}

Scripts~10--13 provide reduced-orientation cross-checks for $S_5$; their
numerical maxima are not uniform bounds over the full ordered universe.
The level-$3$ census in Scripts~$12$--$13$ contains $73{,}810$ retained pairs.
Scripts~$17$ and~$18$ cross-check all $9{,}447{,}840$ orientations numerically.
Script~$20$ supplies the definitive proof-grade certificate: exact symmetry
accounting followed by integer Gram-form inequalities for every relevant
irreducible proves the $71/2000$ bound and the tighter global upper bound
$0.0354002321$. Scripts~$14$ and~$15$ verify
\Cref{thm:dim-obstruction}: Script~$14$ exhaustively enumerates the reduced
level-$2$ and level-$3$ censuses of $S_6$ and $S_7$, reporting
a maximum norm of~$1$ and also strictly contractive configurations. It verifies
explicit norm-$1$ examples and exact Gram-form certificates for contractive
examples in both groups. Script~$15$ gives the exact singular-value pattern
$[\underbrace{1,\ldots,1}_{n-3}, 1/8, 0]$ of $\Dhat_1$ for $S_5, S_6, S_7$.
Script~$16$ verifies the conditioning lemma for $\ell=4$ over all $2^{16}$ paths
in exact rational arithmetic. Scripts~$01$--$09$, $11$--$13$, $17$--$19$
provide additional cross-validation.


\section*{Acknowledgements}

Generative AI tools were used for editing, proofreading, and assistance with
the verification scripts. The author takes full responsibility for the content.


\bibliography{refs}

\end{document}
